\subsection{Multi-output Model}

% Multi-output models
\begin{table}[htbp]
\centering
\caption{Performance of the multi-output models. The best value for each metric is highlighted in bold.}
\label{tab:multi_output_models}
\resizebox{\textwidth}{!}{%
\begin{tabular}{p{5.2cm}cccc}
\toprule
Model &
Dice $\uparrow$ &
RMSE / RMSE$_{\mathrm{whole}}$ $\downarrow$ &
Soft Dice$_{\mathrm{whole}}$ $\uparrow$ &
Recognition $\uparrow$ \\
\midrule
Multi-output Tversky loss, ReduceLROnPlateau factor 0.5
& \textbf{0.9411}
& \textbf{0.0396} / \textbf{0.0084}
& \textbf{0.9758}
& \textbf{0.9507} \\

Multi-output Tversky loss, learning rate $10^{-4}$
& \textbf{0.9411}
& 0.0432 / 0.0104
& 0.9728
& 0.9490 \\

Multi-output Tversky loss, batch size 20
& 0.9380
& 0.0401 / 0.0096
& 0.9729
& 0.9490 \\

Multi-output Tversky loss, learning rate $5\times10^{-5}$
& 0.9349
& 0.0409 / 0.0094
& 0.9686
& 0.9470 \\

Multi-output Tversky loss, batch size 100
& 0.9332
& 0.0411 / 0.0095
& 0.9704
& 0.9460 \\

Multi-output Tversky loss, 32 base channels
& 0.9294
& 0.0431 / 0.0104
& 0.9685
& 0.9432 \\

Multi-output Tversky loss, learning rate $3\times10^{-4}$
& 0.9086
& 0.0492 / 0.0119
& 0.9700
& 0.9297 \\

Multi-output BCE/MSE loss
& 0.8214
& 0.1273 / 0.0309
& 0.7869
& 0.8471 \\

Multi-output BCE/MSE loss, small dataset
& 0.6439
& 0.1460 / 0.0497
& 0.6283
& 0.7490 \\

\bottomrule
\end{tabular}%
}
\end{table}

\begin{figure}[htbp]
    \centering
    \includegraphics[width=0.95\textwidth]{diagrams/multihead/group_04_multi_head_objective_scheduler.pdf}
    \caption[Multi-output objective and scheduler comparison]{Comparison of the multi-output objectives and scheduler settings. The binary cross-entropy/mean squared error baseline is clearly weaker than the Tversky models, while the Tversky model with ReduceLROnPlateau gives the strongest balance of segmentation quality, intensity reconstruction, and spot recognition.}
    \label{fig:multihead-objective-scheduler}
\end{figure}

\subsubsection{Binary Cross-Entropy and Mean Squared Error Loss}

Table~\ref{tab:multi_output_models} shows that the full-dataset \gls{bce}/\gls{mse} model is stronger than the one-output baselines, but remains substantially weaker than the Tversky multi-output models. It reaches Dice $0.8214$, \gls{rmse} $0.1273$ / $0.0309$, soft Dice$_{\mathrm{whole}}$ $0.7869$, and recognition $0.8471$. The selected examples show that this average hides a considerable lower tail: strong \gls{bce}/\gls{mse} cases can be very accurate, while poor cases still fail to recover both separated spots reliably.

The \gls{bce}/\gls{mse} loss was also tested on the smaller dataset. The loss curves in Figure~\ref{appfig:loss-small-dataset} indicate overfitting: the training loss continued to converge, while the validation loss worsened. The model was nevertheless evaluated on unseen validation data. Its metrics are lower than those of the full-dataset \gls{bce}/\gls{mse} model, with Dice decreasing from $0.8214$ to $0.6439$, \gls{rmse} worsening from $0.1273$ / $0.0309$ to $0.1460$ / $0.0497$, soft Dice$_{\mathrm{whole}}$ decreasing from $0.7869$ to $0.6283$, and recognition decreasing from $0.8471$ to $0.7490$. This confirms that the small dataset limits generalization across all metrics shown in the main table. Nevertheless, even this small-dataset multi-output baseline remains competitive with, and in several metrics stronger than, the one-output approaches.

\subsubsection{Tversky Loss and Hyperparameter Tuning}

Switching to the Tversky foreground loss produces the decisive improvement. The selected scheduler configuration in Table~\ref{tab:multi_output_models} reaches Dice $0.9411$, \gls{rmse} $0.0396$ / $0.0084$, soft Dice$_{\mathrm{whole}}$ $0.9758$, and recognition $0.9507$ on the validation data. Relative to the full-dataset \gls{bce}/\gls{mse} baseline, this improves Dice by $0.1197$, foreground \gls{rmse} by $0.0877$, whole-image \gls{rmse} by $0.0225$, soft Dice$_{\mathrm{whole}}$ by $0.1889$, and recognition by $0.1036$.

The hyperparameter-tuning results in Table~\ref{tab:multi_output_models} show that the Tversky loss is robust, but still sensitive to capacity and optimization settings. Reducing the model to $32$ base channels lowers Dice to $0.9294$ and recognition to $0.9432$, and increases \gls{rmse} to $0.0431$ / $0.0104$. The problem example in Figure~\ref{fig:multihead-base32-problem} illustrates the practical effect of this reduction: the smaller model can still solve many samples, but it is more likely to lose structure on difficult overlaps. Increasing the batch size to $100$ also underperforms the best settings, with Dice $0.9332$, \gls{rmse} $0.0411$ / $0.0095$, and recognition $0.9460$. The batch-size-$20$ run is stronger, reaching Dice $0.9380$ and recognition $0.9490$, which suggests that the optimization benefits from a moderately smaller batch.

\begin{figure}[htbp]
    \centering
    \includegraphics[width=0.95\textwidth]{diagrams/multihead/tversky_base_32_problem.pdf}
    \caption[Multi-output base channel problem example]{Problem case for the smaller base channel multi-output Tversky model. The example illustrates that reducing model capacity increases the risk of failures on difficult overlaps.}
    \label{fig:multihead-base32-problem}
\end{figure}

The learning rate tuning shows a similar pattern. A learning rate of $3\times10^{-4}$ is the weakest Tversky setting in the main table, with Dice $0.9086$, \gls{rmse} $0.0492$ / $0.0119$, and recognition $0.9297$. The corresponding loss curves in Figure~\ref{fig:multihead-learning-rate-comparison} show that this run does not converge satisfactorily: the optimization remains unstable compared with the other learning rate settings and does not reach the same validation quality. The lower learning rate of $5\times10^{-5}$ is considerably stronger, but remains below the selected scheduler run. The fixed learning rate of $10^{-4}$ ties the scheduler run on Dice at $0.9411$, but the scheduler gives better \gls{rmse}, soft Dice$_{\mathrm{whole}}$, and recognition. Therefore, the scheduler model was selected as the final model. The learning rate behavior is shown explicitly in Figure~\ref{fig:multihead-tversky-learning-rate}, and the loss components in Figure~\ref{fig:multihead-tversky-validation-components} show that the selected model keeps both mask and intensity terms stable during training.

\begin{figure}[htbp]
    \centering
    \includegraphics[width=0.95\textwidth]{diagrams/multihead/multi_head_tversky_plateau_validation_components.pdf}
    \caption[Multi-output Tversky validation components]{Validation components for the multi-output Tversky model with ReduceLROnPlateau. The component curves show how the final model balances mask and intensity quality.}
    \label{fig:multihead-tversky-validation-components}
\end{figure}

\begin{figure}[htbp]
    \centering
    \includegraphics[width=0.95\textwidth]{diagrams/multihead/multi_head_tversky_plateau_learning_rate.pdf}
    \caption[Multi-output Tversky learning rate schedule]{Learning rate schedule for the multi-output Tversky model with ReduceLROnPlateau. The learning rate is reduced by a factor of $0.5$ after plateaus of $30$ epochs.}
    \label{fig:multihead-tversky-learning-rate}
\end{figure}

\begin{figure}[htbp]
    \centering
    \includegraphics[width=0.95\textwidth]{diagrams/multihead/group_05_multi_head_learning_rate.pdf}
    \caption[Multi-output learning rate comparison]{Comparison of learning rate settings for the multi-output Tversky models. The fixed learning rate runs and the scheduler run show that overly small learning rates underfit, while the scheduler gives the best balance of spot recognition and intensity error.}
    \label{fig:multihead-learning-rate-comparison}
\end{figure}

\begin{figure}[htbp]
    \centering
    \includegraphics[width=0.95\textwidth]{diagrams/multihead/group_06_multi_head_batch_base.pdf}
    \caption[Multi-output batch and base channel comparison]{Comparison of batch-size and base channel settings for the multi-output Tversky models. Reducing the base channels and using a batch size of $100$ both reduce segmentation quality compared with the stronger Tversky settings.}
    \label{fig:multihead-batch-base-comparison}
\end{figure}

Figures~\ref{fig:multihead-learning-rate-comparison} and~\ref{fig:multihead-batch-base-comparison} summarize the hyperparameter-tuning process visually. Together with Table~\ref{tab:multi_output_models}, they show that no single change explains the final performance. The per-dataset heatmap in Figure~\ref{fig:multihead-dataset-heatmap} further shows that the multi-output Tversky models improve the lower-performing material groups compared with the \gls{bce}/\gls{mse} baseline, but that the task remains dataset dependent. Difficult material groups and strongly overlapping examples still form the lower tail, so the aggregate score should be interpreted together with the dataset-wise behavior. The final model is therefore the result of combining the multi-output architecture, the Tversky foreground objective, sufficient model capacity, a moderate batch size, and an adaptive learning rate schedule.

\begin{figure}[htbp]
    \centering
    \includegraphics[width=0.95\textwidth]{diagrams/multihead/example_best_model.pdf}
    \caption[Best multi-output model example]{Representative example for the selected multi-output model. The prediction keeps both mask localization and intensity reconstruction substantially more stable than the one-output and two-output alternatives.}
    \label{fig:multihead-best-example}
\end{figure}

\subsection{Final Model Evaluation on the Test Dataset}

Before applying the final model to the held-out test dataset, the most relevant validation candidates were compared across the output formulations. Table~\ref{tab:final_model_candidate_comparison} summarizes this selection step. The one-output Dice/L1 model remains physically limited by its residual construction, and the two-output L1 model with batch size $100$ is strong but still couples localization and intensity reconstruction in the same output image. The multi-output \gls{bce}/\gls{mse} baseline improves on the one-output formulation, but it is clearly weaker than the multi-output Tversky scheduler model. The ReduceLROnPlateau configuration is therefore selected as the final model because it reaches the best \gls{rmse}, best soft Dice$_{\mathrm{whole}}$, and best recognition while also achieving the highest Dice together with the fixed $10^{-4}$ run.

\begin{table}[htbp]
\centering
\caption{Validation comparison of the most relevant model candidates.}
\label{tab:final_model_candidate_comparison}
\begin{tabular}{p{5.2cm}cccc}
\toprule
Model &
Dice $\uparrow$ &
RMSE / RMSE$_{\mathrm{whole}}$ $\downarrow$ &
Soft Dice$_{\mathrm{whole}}$ $\uparrow$ &
Recognition $\uparrow$ \\
\midrule
One-output Dice/L1 loss
& 0.6843
& 0.1350 / 0.0337
& 0.7997
& 0.7747 \\

Two-output L1 loss, batch size 100
& 0.9222
& 0.0482 / 0.0096
& 0.9419
& 0.9370 \\

Multi-output BCE/MSE loss
& 0.8214
& 0.1273 / 0.0309
& 0.7869
& 0.8471 \\

Multi-output Tversky loss, ReduceLROnPlateau factor 0.5
& \textbf{0.9411}
& \textbf{0.0396} / \textbf{0.0084}
& \textbf{0.9758}
& \textbf{0.9507} \\
\bottomrule
\end{tabular}
\end{table}

The corresponding per-dataset heatmaps in Figures~\ref{fig:oneoutput-dataset-heatmap},~\ref{fig:twooutput-per-dataset}, and~\ref{fig:multihead-dataset-heatmap} show that the candidate ranking is not uniform across all source datasets. Across the compared models, the following dataset-level tendencies are visible:
\begin{itemize}
    \item \texttt{IN718\_twins} is the most consistently favourable dataset. It contributes many of the strongest one-output and two-output examples and remains well handled by the multi-output models, suggesting that compact, well-localized spot morphologies are easier to separate.
    \item The aluminium datasets with smaller or regular grains are generally more stable than the difficult material groups. Their performance improves strongly with the multi-output Tversky model, indicating that the final architecture can use the additional mask and intensity heads to reduce dataset-specific failures.
    \item \texttt{Ti7Al}, \texttt{Iron}, \texttt{Iron\_deformed}, and partly \texttt{Al\_big\_grains} form the lower-performing group in the one-output and two-output comparisons. These datasets contain more difficult morphology or overlap geometry, so the models more often recover only one spot or produce less accurate intensity separation.
    \item An interesting pattern is that \texttt{IN718\_twins} appears in both strong and weak examples. Dataset identity alone therefore does not explain the result; the overlap geometry and individual spot morphology still matter within a nominally favourable dataset.
    \item The multi-output Tversky model reduces the gap between easy and difficult datasets most clearly. The remaining weak cases are no longer a failure of one specific material only, but are concentrated in high-overlap or irregular samples across the harder groups.
\end{itemize}

\begin{table}[htbp]
\centering
\caption{Performance of the selected final model on the held-out test dataset.}
\label{tab:final_model_test_evaluation}
\begin{tabular}{p{5.2cm}cccc}
\toprule
Model &
Dice $\uparrow$ &
RMSE / RMSE$_{\mathrm{whole}}$ $\downarrow$ &
Soft Dice$_{\mathrm{whole}}$ $\uparrow$ &
Recognition $\uparrow$ \\
\midrule
Multi-output Tversky loss, ReduceLROnPlateau factor 0.5 evaluation
& 0.9300
& 0.0460 / 0.0105
& 0.9581
& 0.9398 \\
\bottomrule
\end{tabular}
\end{table}

The selected model was then evaluated on a held-out test dataset of $20{,}000$ synthetically generated overlaps. Table~\ref{tab:final_model_test_evaluation} reports Dice $0.9300$, \gls{rmse} $0.0460$ / $0.0105$, soft Dice$_{\mathrm{whole}}$ $0.9581$, and recognition $0.9398$. These values are lower than the validation metrics of the selected scheduler run, but they remain close enough to indicate that the model generalizes well to unseen generated samples. The representative prediction in Figure~\ref{fig:multihead-best-example} shows the main advantage of the final model: the two masks remain spatially separated while the corresponding intensity heads preserve the spot signal.

\begin{figure}[htbp]
    \centering
    \includegraphics[width=0.95\textwidth]{diagrams/multihead/figure_04_dataset_model_heatmap_multi_output.pdf}
    \caption[Multi-output per-dataset metrics]{Per-dataset metric heatmap for the multi-output models. The heatmap shows that the multi-output Tversky models improve performance across datasets, while difficult material datasets still form the lower tail.}
    \label{fig:multihead-dataset-heatmap}
\end{figure}

\begin{figure}[htbp]
    \centering
    \includegraphics[width=0.95\textwidth]{diagrams/multihead/figure_05_spot_counter_multi_output.pdf}
    \caption[Multi-output spot-count performance]{Spot-count grouped performance for the multi-output models. The final Tversky scheduler model recovers both spots more consistently than the binary cross-entropy/mean squared error baseline and the reduced-capacity models.}
    \label{fig:multihead-spot-counter}
\end{figure}

The dataset heatmap in Figure~\ref{fig:multihead-dataset-heatmap} and the spot-count plot in Figure~\ref{fig:multihead-spot-counter} show that the final model is not perfect, but its errors are much less severe than in the previous output models. The best selected cases reach very high recognition and Dice values with very low \gls{rmse}, while the remaining weak cases are concentrated in high-overlap or difficult material examples. Nevertheless, the multi-output Tversky model with ReduceLROnPlateau is the strongest evaluated model because it combines accurate foreground separation with substantially improved intensity reconstruction and a much smaller failure tail.
